\section{Introduction}

Suppose an abelian variety has good reduction at a prime. A rational
algebraic cycle on the variety specializes to a cycle on the reduction.
Its intersection with any complementary product of divisors has a rational
degree, and every cycle-class realization computes that same number.
Milne's rationality conjecture asks for this conclusion when the original
class is only known to be a rational Hodge class. We prove it for every
residue characteristic.

The distinction matters because good reduction can acquire divisors that
do not lift to the original variety. Deligne's absolute-Hodge theorem gives
compatible realizations of the Hodge class, but does not by itself identify
its pairing with these new divisors as the degree of a rational cycle.
The conjecture isolates this rationality question from the stronger problem
of representing the Hodge class by an algebraic cycle.

Fix \(\Qbar\subset\C\), a prime \(p\), and a \(p\)-adic place \(w\)
of \(\Qbar\). Write
\[
 \F=\overline{\mathbb F}_p,\qquad
 \Kzero=\operatorname{Frac}W(\F),\qquad
 \Cw=\widehat{\Qbar}_w.
\]
We regard \(\Kzero\) and an algebraic closure of it as subfields of
\(\Cw\). Good reduction means good reduction after descent to a number
field and finite extension at the chosen place.

\begin{theorem}\label{thm:main}
Let \(A/\Qbar\) be an abelian variety of dimension \(d\) with good
reduction \(A_0/\F\) at \(w\). Let \(0\leq r\leq d\), and let
\[
 \gamma\in H^{2r}(A(\C),\Q(r))
\]
be a rational Hodge class. For every list of Cartier divisors
\(D_1,\ldots,D_{d-r}\) on \(A_0\), there is a number \(q\in\Q\)
such that
\begin{align*}
 \Tr_\ell\bigl(\gamma_{0,\ell}
       c_{1,\ell}(D_1)\cdots c_{1,\ell}(D_{d-r})\bigr)&=q
       &&\text{in }\Q_\ell\quad(\ell\ne p),\\
 \Tr_p\bigl(\gamma_{0,p}
       c_{1,\mathrm{cris}}(D_1)\cdots
       c_{1,\mathrm{cris}}(D_{d-r})\bigr)&=q
       &&\text{in }\Cw.
\end{align*}
Here \(\gamma_{0,\ell}\) is the smooth proper specialization of the
\(\ell\)-adic realization, and \(\gamma_{0,p}\) is the
de Rham--crystalline specialization in
\[
 H^{2r}_{\mathrm{cris}}(A_0/W(\F))[1/p](r)
       \otimes_{\Kzero}\Cw.
\]
The traces use the top Tate twist and send the class of a zero-cycle to
its degree. If \(r=d\), the divisor product is \(1\).
\end{theorem}

By linearity the same conclusion holds for every rational Lefschetz class
of complementary codimension, that is, every rational linear combination
of products of divisors. The number \(q\) is common to all realizations
for the fixed data; no comparison between different reduction places is
required.

\subsection{Origins and consequences}

Milne introduced a weak-rationality formulation while studying the
reduction of Shimura varieties and the Langlands--Rapoport conjecture
\cite[Conjecture~6.1(A)]{Milne2000}. It asks that a specialized CM
Hodge class whose components lie in the local spans of divisor products
in every realization arise from one rational Lefschetz class.
The pairing formulation appears in his
AIM notes and his study of rational Tate classes
\cite[Section~10]{MilneAIM2007}, \cite[Conjecture~4.1]{Milne2009}.
A previous positive case is a CM abelian variety with simple ordinary
reduction, together with all its powers: in that case every divisor
class on the reduction lifts \cite[Example~4.2]{Milne2009}.

The broader purpose is to obtain rational cohomological correspondences
for abelian motives in characteristic~\(p\), compatible with reduction
from characteristic zero \cite[Sections~5--6]{Milne2000}; see also
\cite[Introduction and Section~2]{Milne2025}. For \(\ell\ne p\),
Tate classes are twisted \(\ell\)-adic classes fixed by a power of
Frobenius; there is a corresponding \(\Q_p\)-algebra of crystalline
Tate classes. A \emph{good theory of rational Tate classes} chooses
one graded \(\Q\)-algebra on each variety whose scalar extensions
give these Tate algebras. The theory is stable under pullback and
pushforward, contains divisor classes, and contains specializations of
CM Hodge classes \cite[Definitions~2.1 and~3.1]{Milne2009}.
These are cohomological correspondences; their definition does not
assert that they are algebraic.

Milne proved that the rationality assertion for all CM abelian varieties
is equivalent to the existence of such a good theory on varieties over
\(\F\) whose connected components are products of abelian varieties
and projective spaces \cite[Theorem~4.5]{Milne2009}.
Our proof also expresses each specialized Hodge class as the pullback,
along a rational homomorphism of special fibers, of a specialized Hodge
class from a CM abelian variety. This extends the theory's
specialization property beyond the CM case.
In Corollary~\ref{cor:all-hodge-specialization}, every rational
Hodge class on a good-reduction abelian variety specializes to a unique
class in one fixed good theory. Its crystalline component gives the
stated specialization after extension to \(\Cw\).

A distinct additional input gives algebraicity on the special fiber.
Combining the same transport with the Hodge theorem for CM abelian
varieties \cite[Theorem~1.1]{OpenAICMHodge2026},
Corollary~\ref{cor:algebraic-specialization} represents each specialized
Hodge class by one rational algebraic cycle in all the stated
realizations. Its pairing with any complementary algebraic cycle is
therefore a rational intersection number. This consequence uses the CM
Hodge theorem; the proof of Theorem~\ref{thm:main} is independent of it.

\subsection{The method and its antecedents}

Three foundational results make this question accessible. Deligne proved
that every Hodge class on an abelian variety is absolute Hodge
\cite[Main Theorem~2.11]{Deligne1982}. Blasius established its compatibility
with \(p\)-adic comparison \cite[Theorems~0.3 and~5.3]{Blasius1994};
Moonen presents the Blasius--Wintenberger comparison theorem, with
Ogus's simplification, in the perfect-residue-field form needed for
the local auxiliary lifts below \cite[Theorem~5.6.3]{Moonen1998}. Milne's theory of
Lefschetz classes supplies a common rational divisor algebra with perfect
complementary pairing, independent of the Weil cohomology
\cite[Section~5]{Milne1999}.

The determinant classes used in the proof belong to the theory of Weil
classes developed by Weil, Deligne and Andr\'e; see
\cite[Section~4 and Endnotes~M.10--M.12]{Deligne1982}.
Andr\'e's refinement of the CM reduction expresses Hodge classes by
pullbacks from varieties of split Weil type \cite{Andre1992}; a detailed
account is \cite[Theorem~1]{Milne2020}. We prove directly the balanced
form of that reduction needed here in Lemma~\ref{lem:cm-reduction}.
Balanced signatures assert equal Hodge multiplicities at conjugate
embeddings; split Weil type also requires the associated Hermitian
form to have a half-dimensional isotropic subspace.
Andr\'e's later theory of motivated cycles establishes a broader
unconditional motivic setting for abelian Hodge classes
\cite[Theorem~0.6.2]{Andre1996}, without identifying them with algebraic
cycles.

For an imaginary quadratic field, Markman has proved algebraicity of
Weil classes on abelian fourfolds of Weil type and on abelian sixfolds
of split Weil type \cite[Theorem~1.2]{Markman2026}. These are
characteristic-zero algebraicity results. The pairing argument here
applies in arbitrary dimension and begins only with the Hodge class
and its good reduction.

There is also a direct antecedent for the lifting strategy. Milne asked
whether a specialized CM Hodge class and a divisor could be lifted
together, up to isogeny, to a CM variety, and explained the resulting
rationality implication \cite[Question~10.6 and Proposition~10.7]{MilneAIM2007}.
His current account gives a conditional criterion through weak lifting
of Lefschetz classes on CM Weil triples and describes the reduction
from general Hodge classes through CM varieties to Weil classes
\cite[Introduction, Question~2.17, and Theorem~2.18]{Milne2025}.
Our proof follows this general route, using two auxiliary varieties
with different purposes. The
first is CM; the second realizes a prescribed admissible filtration
over a local field and need not be CM. The necessary realization
machinery comes from Grothendieck--Serre--Tate deformation theory,
Tate's extension theorem, and the integral-model and local-Shimura
results of Kisin--Zhou and Pappas--Rapoport
\cite{Grothendieck1970,Tate1967,KisinZhou2025,PappasRapoport2024}.
For the prescribed filtration, we choose a framed local Shimura
point over its admissible period and use integral uniformization;
the required uniformization condition is supplied by
\cite[Corollary~6.3]{GleasonLimXu2026}. We state the precise inputs
where they enter and prove the connecting arguments for the
prescribed flag.

The contribution along this route is to construct a filtration that
preserves the extra divisor and then obtain a rational divisor identity
on its lift. A separate normalization by rational intersection numbers
makes the tensor transports simultaneous in all realizations.
Kisin--Zhou's compatibility theorem concerns Weil--Deligne
representations across coefficient fields
\cite[Theorem~1.2]{KisinZhou2025}; the additional issue here is the
rational value of a Hodge--divisor intersection.

The key step is a dichotomy for balanced Weil classes: either all the
complementary Lefschetz pairings vanish, or the specialized class is
itself a rational polynomial in divisors, simultaneously in every
realization.

\subsection{Proof overview}

The proof has four stages. The order below describes the deduction;
Sections~\ref{sec:transport}--\ref{sec:weil} establish its reusable
constructions before Section~\ref{sec:reduction} assembles them.

\paragraph{1. Transport to a CM variety.}
Corollary~\ref{cor:cm-transport} constructs a CM
abelian variety \(C/\Qbar\), a quasi-isogeny
\(a_{\mathrm{CM}}:A_0^n\to C_0\), and a rational Hodge class
\(\kappa\) such that
\[
 \gamma_0=(a_{\mathrm{CM}}j)^*\kappa_0,
\]
where \(j\) includes the first factor. This is one rational algebraic homomorphism on the
special fibers and one equality in every realization. Kisin--Zhou's
special-point theorem supplies the CM lift
\cite[Theorem~2.2.7]{KisinZhou2025}. The local tensors initially determine
similitude scalars; a ratio of rational divisor intersections shows
that all those scalars are the same rational number. This normalization
is what transports the original rational class, rather than a separate
local multiple in each cohomology theory.

\paragraph{2. Reduce CM classes to balanced Weil classes.}
Lemma~\ref{lem:cm-reduction} constructs rational-idempotent factors of
actual powers of \(C\). Their homology has rank \(2s\) over a CM field
\(E\), with signatures \((s,s)\) at every embedding, and their
polarizations induce conjugation on \(E\). Their Weil classes are the
rational forms whose components, after splitting \(E\), are top
alternating forms on individual \(E\)-eigenspaces. Pullbacks of these
Hodge classes span the given CM Hodge class over \(\Q\). Good reduction
is retained. Having one rational value for each complementary
Lefschetz pairing is preserved by rational pullback
(Lemma~\ref{lem:pairing-functorial}). It therefore remains to prove
that assertion for balanced Weil classes.

\paragraph{3. Find a divisor and a compatible filtration.}
For a balanced Weil variety \(Z\), let \(\dagger\) denote Rosati on its
reduction and consider
\[
 \mathcal U=\{u\in\End^0(Z_0):u^\dagger=u,\quad
                         ue=\bar e u\text{ for all }e\in E\}.
\]
Proposition~\ref{prop:weil} proves a dichotomy. If \(\mathcal U\) has no
invertible element, every complementary Lefschetz pairing of a
specialized Weil class is zero. If it has an invertible element, the
specialized Weil class is itself a simultaneous rational Lefschetz
class: one rational polynomial in divisors gives its class in every
realization. For the second branch, such a \(u\) supplies a divisor
\(D^u\) whose alternating form is nondegenerate on each \(E\)-eigenspace.
Theorem~\ref{thm:filtration} constructs a weakly admissible balanced
filtration on the fixed crystalline module that preserves \(E\), \(u\),
and the polarization. Generic transversality has a classical role in
admissible-filtration constructions
\cite[Theorem~1 and Lemmas~1.1--1.2]{FontaineRapoport2005}.
Here bounded-degree avoidance makes all the constrained choices over
one finite extension, and a maximality argument reduces the
weak-admissibility test to subobjects stable under the endomorphisms.

\paragraph{4. Lift the divisor and extract a rational identity.}
The second case of Proposition~\ref{prop:transport} realizes this flag
on a local auxiliary abelian variety \(B\), with a marked quasi-isogeny
\(a_{\mathrm{fl}}:Z_0^m\to B_0\). The prescribed flag makes the
first-summand projector, its algebraic \(E\)-action, and the transported
divisor filtered. These specified algebraic maps lift to \(B\).
After choosing an abstract complex realization of \(B\), polynomial
coefficient extraction expresses the relevant rational Betti Weil class
as a rational linear combination of divisor products
(Lemma~\ref{lem:weil-divisor-span}). The same rational identity specializes
in every realization. Pullback returns it to \(Z_0\), completing the
case of an invertible element in the dichotomy and, through the first two stages, the
main theorem.

Figure~\ref{fig:two-transports} separates the two lifts and their
special-fiber markings. Section~\ref{sec:preliminaries} fixes comparisons,
traces, and rational pullback. Section~\ref{sec:transport} proves both
transport alternatives; Section~\ref{sec:filtration} constructs the
prescribed flag; Section~\ref{sec:weil} proves the balanced-Weil dichotomy;
and Section~\ref{sec:reduction} proves the CM reduction and finishes the
argument.

\begin{figure}[tbp]
\centering
\begin{tikzpicture}[
 >=Latex, font=\small,
 variety/.style={draw=MidnightBlue!60,rounded corners=2pt,
                 fill=MidnightBlue!3,align=center,minimum width=43mm,
                 minimum height=11mm},
 every edge/.style={draw,->,thick},
 note/.style={font=\footnotesize,align=center}]
\node[font=\bfseries] at (3,1.0) {CM transport};
\node[variety] (A) at (0,0) {$A^n/\Qbar$};
\node[variety] (C) at (6,0) {$C/\Qbar$\quad(CM)};
\node[variety] (A0) at (0,-1.7) {$A_0^n$};
\node[variety] (C0) at (6,-1.7) {$C_0$};
\path (A) edge node[left,note] {reduction} (A0);
\path (C) edge node[right,note] {reduction} (C0);
\path (A0) edge node[above] {$a_{\mathrm{CM}}$} (C0);
\node[note] at (3,-2.65) {$\gamma_0=(a_{\mathrm{CM}}j)^*\kappa_0$};
\node[font=\bfseries] at (3,-3.3) {Prescribed-filtration transport};
\node[variety] (Z) at (0,-4.4) {$Z^m/\Qbar$\quad(balanced Weil)};
\node[variety] (B) at (6,-4.4) {$B/K$, $K/K_0$ finite};
\node[variety] (Z0) at (0,-6.1) {$Z_0^m$};
\node[variety] (B0) at (6,-6.1) {$B_0$};
\path (Z) edge node[left,note] {reduction} (Z0);
\path (B) edge node[right,note] {reduction} (B0);
\path (Z0) edge node[above] {$a_{\mathrm{fl}}$} (B0);
\node[note] at (3,-7.05)
 {$a_{\mathrm{fl}}^*\Fil^1H^1_{\mathrm{dR}}(B)
       =(\mathcal F_Z^1)^{\oplus m}$};
\end{tikzpicture}

\caption{The two uses of transport. Each lower horizontal arrow is one
actual quasi-isogeny preserving the specified tensors in all
realizations; the vertical arrows denote good reduction. There is no
claimed map between the generic fibers. The upper construction starts
from the original variety. The lower construction is applied separately
to a balanced Weil variety \(Z\) arising in the CM reduction; its lift
\(B\) realizes the chosen filtration \(\mathcal F_Z^{\oplus m}\) under
\(a_{\mathrm{fl}}^*\). The CM and prescribed-filtration auxiliaries need
not coincide.}
\label{fig:two-transports}
\end{figure}
